% New finite recursion, 2026-10-08. No imported old H/h tables.
\subsection{A finite recursion on typed stable-prefix states}
\label{sec:new-finite-double-sieve-recursion}

This construction uses only the explicitly proved first and replacement
second seeds, their stable-prefix extensions, and finite Buchstab
identities. It does not transplant any old correction table, claim an
optimal correction function, or take an infinite iteration limit.
Its numerical implementation is
\texttt{probe\_new\_finite\_double\_sieve.py}; the saved JSON is a
floating feasibility calculation, not an outward certificate.

At a state retain the marked specified prefix $B=N^{T+o(1)}$, one
total level $Q=N^\nu$, residual $R=Q/B=N^{L+o(1)}$, and
$\alpha=(1-T)/L\ge2$. The state also retains its enhanced
factorization branch, all its gates and square tests, and a ceiling
$\chi$ below which that branch persists. For a $W$ branch with a
future third-slot bound $c$, require $\chi\le c$ as well as
$\chi$ below the smallest specified slot. Alternatively the state
uses the classical absolute Bombieri--Vinogradov interface at
$\nu<1/2$. All physical cutoffs used below are below the retained
ceiling. Enhanced non-switched nodes retain full well support;
positive switched nodes use the separate finite-product theorem.

Put $A(s)=sF(s)/(2e^\gamma)$ and $a(s)=sf(s)/(2e^\gamma)$.
For fixed generation $n$, the claims represented by $H_n,h_n$ are
the simultaneous bounds
\[
 \Phi_s\le\{A(s)-H_n(s)+o(1)\}\Theta_Q,\qquad
 \Phi_s\ge\{a(s)+h_n(s)+o(1)\}\Theta_Q,
\]
uniformly over the indicated typed states and every $\alpha\ge2$.
Only upper nodes at $s\ge1$ and lower nodes at $s\ge2$ are used;
set $h_n(s)=0$ below two. The initial seed is
\[
 H_0(s)=\frac1{50}\mathbf1_{1\le s\le11/5},\qquad h_0(s)=0.
\]
Its strict uniform inequality is separately certified in
\texttt{second-seed-omega-certified.tex}. Zero is chosen outside
the displayed support. A finite positive prefix partition licenses
all choices of upper or lower estimates.

\paragraph{A common upper envelope for the Buchstab function.}
Let $C=70893/125000=.567144$. The exact inequality
$C\sum_{j=0}^{8}C^j/j!>1$ implies $Ce^C>1$, hence
$C+\log C>0$. On $2\le u\le3$ the maximum of
$1+\log(u-1)-Cu$ is $-\log C-C<0$; the delay equation propagates
$\omega(u)\le C$ for $u\ge3$. Since $\omega(u)=1/u$ on $[1,2]$,
\[
 \omega(v)\le\Omega(u):=\max\{1/u,C\}\qquad(v\ge u\ge1).
\]
Consequently the first-seed switched integral is bounded uniformly
for all $\alpha\ge2$ by
\[
 \overline I(s,s_0)=
 \int_{1/s_0<u_1<u_2<u_3<1/s}
 \Omega\!\left(\frac{2-u_1-u_2-u_3}{u_2}\right)
              \frac{du_1du_2du_3}{u_1u_2^2u_3},\qquad s\ge2.
\]
Its argument is at least $2s-3\ge1$. This is a pointwise tail
majorant, not an assumption that the switched integral attains its
supremum at $\alpha=2$. Unit quotients in this triple term are
power-saving; the second seed retains its separate five- and
six-prime unit masses.

\paragraph{The positive finite operators.}
Fix $M=10$, and choose zero defects at any unproved argument beyond
the finite table. For $s\ge2$, the following are valid new candidates:
\begin{align*}
 \mathcal B_H[h_n](s)&=\int_{\max(2,s-1)}^M h_n(w)\frac{dw}{w},\\
 \mathcal B_h[H_n](s)&=\int_{s-1}^M H_n(w)\frac{dw}{w},
\end{align*}
where a reversed interval contributes zero. These identities start
the ordinary comparison at argument $M+1$ and discard any correction
there. Ordinary baseline cancellation is exact at that common level.
In particular upper children with physical argument $1\le w<2$
are included in the lower operator. Their defect comes from a
smaller comparison cutoff, while the lower function below two is
still zero.

The first weighted inequality gives the further upper candidate
\begin{align*}
 \mathcal W_1[H_n,h_n](s;s_0)
 ={}& A(s)-A(s_0)+H_n(s_0)-\frac K8\overline I(s,s_0)\\
 &+\frac12\int_{1/s_0}^{1/s}
       \frac{a(s_0(1-u))+h_n(s_0(1-u))}{u(1-u)}\,du,
 \qquad 2\le s<s_0\le10.
\end{align*}
The feedback is evaluated at the same residual levels as the
ordinary terms. It is not added at an old argument.
The earlier first-seed search used $s_0\le5$. Wu's pointwise
twofold inequality also holds for every fixed $2\le s<s_0\le10$.
In that larger range use the full $A,a$, not their initial-interval
logarithmic simplifications. The negative-term prime still satisfies
$q^2\le R$, the triple unit-quotient gap remains
$\alpha-3/s\ge1/2$, and every switched prime has exponent at least
$L/s_0>0$. Thus the same signed-prefix and fixed-positive-cutoff
arguments prove exactly the extension used here.

For each admissible finite parameter tuple of the second seed, the
operator in \texttt{second-seed-finite-feedback.tex} contributes
\[
 \mathcal W_2[H_n,h_n](s)=J_2(s)+\frac15
 \left\{4H_n(s_0)+H_n(\phi_1)
       +\sum_{i=2}^4\int h_n(\rho_i)\,d\lambda_i
       +\sum_{i=5}^8\int_{\mathcal G_i}H_n(\rho_i)\,d\lambda_i
 \right\}.
\]
Here $\mathcal G_5$ is only the good square-supported part of term 5;
the bad part remains a positive switched cost. For terms 6--8 it is
the full verified domain. The coefficient measures and child
arguments are precisely those derived in that separate fragment.
The finite parameter bank is recorded by the probe; the chosen tuple
for every generation and target argument is saved in its history.

Take the maximum of the previous defect, zero, and these separately
valid candidates. Then apply positive-count monotonicity to the
upper bound, in its full normalization:
\[
 H_{n+1}(s)\ge
 \sup_{\max(2,s)\le v\le M}
       \{A(s)-A(v)+\widetilde H_{n+1}(v)\}.
\]
For $s\ge2$ retain the candidate at $v=s$ as well. The shift
$A(s)-A(v)$ is essential when an argument exceeds three. No lower
monotonicity with an enlarged physical cutoff is invoked. Maxima
select alternative estimates for a count; overlapping defects are
never added.

\paragraph{Why every generation stays in the typed family.}
A Buchstab child has $q=R^{1/(w+1)}$, $w\ge s-1\ge1$.
Thus $q^2\le R$, $q$ is below the parent's physical cutoff, and
its updated ratio is
\[
 \alpha' =\alpha+\frac{\alpha-1}{w}\ge2.
\]
The complete square condition and selected-remaining conditions
follow by multiplying the residual test by the old prefix.
The new cutoff is at most $q$, which is also below the old stable
ceiling. If $w<2$, its upper comparison is made at $v\ge2$ and
has cutoff $(R/q)^{1/v}\le q$. The subsequent square test is
$q(q')^2\le q(R/q)=R$; it does not require $3\log_Nq\le L$.

The same tests were checked separately for each first-seed and
second-seed non-switched term. Their cap data persist below the
current ceiling, and their ratio after adding exponents remains
at least two. Consequently the functions may be taken uniform in
$\alpha$, even though the actual ratio increases along the tree.
At any fixed finite depth the residual exponents and all newly
specified prime exponents have a fixed positive lower bound on
compact boxes. It can depend on the depth. The finite-product $K$
proof rechooses its Heath--Brown precision for that bound; no
uniformity in infinite depth or in a cutoff tending to zero is
claimed. The modulus coefficient $K$ and its density integral are
independent of this fixed factor-count bound.
For the newer fourth-prefix and tail modulus weights this extension
uses the same original-group argument as
\texttt{generalized-one-step-double-sieve.tex}: retain the original
Heath--Brown scales, choose $12\zeta<\epsilon/2$, and keep the number
of groups fixed before $N\to\infty$. Those new modulus weights
already have the normalized factorization and local-scale margins
needed by the same Type II and long-variable inputs. The coprime
mean and prime-power errors still save a fixed power. Changing this
fixed prime-factor cutoff therefore leaves the modulus-weight density
and its numerical $K$ unchanged. This is an extension of the positive
count theorem, not a signed switched-distribution assertion.

\paragraph{Insertion in the full Chen functional.}
Use the free-threshold eleven-term identity with the complete
$Y_{11}$ range. The probe uses $K=1697/250=6.788$, from the separately
certified fourth-prefix tail refinement, in the five original
switching terms and the positive costs of the two seeds.
The three-parameter range is the positive-$a$ extension of the
earlier range, retaining all its other strict inequalities.
The unrestricted $L_1$ term keeps its established ordinary baseline;
no unproved defect for its different prefix family is inserted.

For $L_2$ add the upper defect in its exact least-prime identity.
For $U_3,U_4$ replace the normalized ordinary upper multiplier by
$A(s)-H_n(s)$ at the actual selected level. For $L_5,L_6$ replace
$a(s)$ by $a(s)+h_n(s)$ only where their full square tests and
$\alpha\ge2$ hold. The latter gate matters for small $a$ on the
constant branch of $V_{2,a}$; the probe retains ordinary lower
bounds and chooses zero extra defect if that gate fails.

On each eligible four-prefix box use
\[
 \min\left\{\frac Kz\omega\left(\frac{1-T}{z}\right),
             \frac4{\nu-T}
               \left[A\left(\max\{2,(\nu-T)/w\}\right)
                     -H_n\left(\max\{2,(\nu-T)/w\}\right)\right]
       \right\}.
\]
All old cap and gate tests are retained. The full $A$ is required
when $a<1/12$; it is not replaced by one. This minimum is evaluated
directly in the functional, so neither its ordinary improvement nor
its switched improvement is counted twice. Each choice is made on
a common finite positive prefix partition before taking limits.

The implementation uses finite arrays, interpolation, Gauss
quadrature and local parameter searches to diagnose the resulting
explicit operators. Those approximations are not a proof that the
computed array entries are outward lower bounds, even though the
underlying exact finite operators preserve the stated inequalities.
An observed small generation change is not an infinite-iteration
error bound or an optimality certificate.

\paragraph{Numerical decision and a hypothetical coefficient test.}
With 32 generations, $\Delta s=.01$, $\Delta s_0=.05$, and ten
Gauss nodes in each triple coordinate, the optimized diagnostic tuple
is approximately
\[
 (a,h,d)=(.08315180,.15198811,.23634369).
\]
The complete functional at quadrature orders 16, 24 and 36 is
respectively $1.93228169$, $1.93234623$, and $1.93235164$.
At $K=6.788$ its zero-margin imbalance diagnostic is $.75226117$;
the remaining coefficient gap to the internal comparison $b=.749$
is approximately $.02949097$. At the original sampled tuple the
generation-24 to generation-32 increase is only $.00002118$.
This is evidence about this finite construction, not an upper bound
on all possible double sieves.

For a decision about future prefix work, keep this exact same
floating $H,h$ array and threshold tuple, and hypothetically lower
only the coefficient of the five original switched counts and the
bad-semiprime removal. Re-evaluating the full four-prefix minimum
gives
\[
\begin{array}{c|c|c}
 \text{assumed }K&c\text{ diagnostic}&c-K\log(1/.749)\\\hline
 6.786&1.93313658&-.02812800\\
 6.750&1.94726640&-.00359360\\
 6.745&1.94922902&-.00018589\\
 6.740&1.95119168& .00322185\\
 6.735&1.95315437& .00662962
\end{array}
\]
The crossing is near $K=6.74473$. In particular a future valid
$K=6.735$ would make this sampled finite construction a plausible
route past the internal comparison. That $K$ is hypothetical here;
the table is not a proof of it, a certified functional value, or a
new imbalance theorem. No reoptimization or further recursive gain
from lowering $K$ is included in this sensitivity test.
